\section{Evaluation Setup}
\label{sec:setup}

To assess \hesp{}, we structure our evaluation around four questions:
\begin{enumerate}
  \item \emph{Reliability:} Does \hesp{} make small local models reliable triage agents, and does it depend on oracle knowledge of the environment? (\S\ref{sec:rq-help})
  \item \emph{Ranking versus Choice:} Is the gain due to the information-gain ranking, or merely to taking the choice of probe away from the model? (\S\ref{sec:rq-rank})
  \item \emph{Probing versus Stopping:} Are choosing what to probe and deciding when to stop separate failures, and which one does each model have? (\S\ref{sec:rq-stop})
  \item \emph{Failure Modes:} How do verdicts fail, in security terms? (\S\ref{sec:security})
\end{enumerate}
Each question was the subject of a pre-registered study designed after the previous one had run (\cref{tab:studies}). \S\ref{sec:realdata} then asks whether public security data could test the same questions.

\begin{table*}[t]
  \caption{Overview of the four pre-registered studies on sec-triage, 7{,}272 episodes in total. The repository labels the studies by version.}
  \label{tab:studies}
  \centering\small
  \begin{adjustbox}{max width=\textwidth}
  \begin{tabular}{llllrl}
    \toprule
    Study & Version & Question & Models & Episodes & Primary endpoint(s) \\
    \midrule
    Initial & v0.5 & does the controller help? & 3 Qwen & 1{,}944 & full \hesp{} vs.\ Memory-only, per model \\
    Table-source & v0.6 & must the tables come from an oracle? & 3 Qwen & 1{,}728 & counted-table \hesp{} vs.\ Memory-only, 7B \\
    Decomposition & v0.8 & ranking or loss of choice? second family & 5 & 1{,}800 & ranking vs.\ random (Q-7B); \hesp{} vs.\ Memory-only (L-8B) \\
    Stopping & v0.9 & what to probe vs.\ when to stop & 5 & 1{,}800 & controller stop; selection given stop (both L-8B) \\
    \bottomrule
  \end{tabular}
  \end{adjustbox}
\end{table*}

\subsection{Environment and Tables}

\subsubsection{The sec-triage Environment} Each episode is one alert on a web service. The hidden cause is one of eight explanations: three attacks (credential stuffing, SQL-injection probing, DNS command-and-control), an insider exfiltrating data, an exposed admin panel, a known-vulnerable component, an authorized vulnerability scan, and a monitoring false positive. The hypothesis set adds \textit{other}. Ten read-only probes, costing 1 to 3 units each, cover authentication failures, HTTP patterns, source reputation, admin exposure, per-user egress, DNS characteristics, component inventory, change tickets, threat intelligence, and a generic runbook. We designed the environment so that shortcuts fail. Threat intelligence says only that an indicator is malicious, and three attacks share that answer. The runbook carries no information. Several benign causes look like attacks until the right probe runs. Each cause appears in three variants. In \textit{base}, outcomes are deterministic except for a 10\% threat-intelligence lag. \textit{Noise} adds 15\% transient HTTP 503 failures and a 30\% lag. In \textit{drift}, the incident is superseded after the second probe: the state version advances and the true cause switches. This gives 24 tasks, and each study runs every task three times per configuration. Each episode is served from its own loopback HTTP sandbox.

\subsubsection{Prediction Tables} We compare three sources of $P(o \mid h, a)$. \emph{Designer} tables come from the environment's generating function and serve as an oracle reference. \emph{Counted} tables (``empirical'' in the tables) are estimated as in \S\ref{sec:tables}, with $k \in \{1, 5, 20, 100\}$ development episodes per cause and variant, from the base and noise variants only and with seeds drawn from a space disjoint from evaluation. \emph{Elicited} tables are produced by the planner LLM itself. Every table is frozen and hashed before any evaluation episode runs. Because development and evaluation episodes come from the same generator, the counted tables tell us how many observations are enough, not how the method behaves in an unfamiliar environment.

\begin{figure*}[t]
  \centering
  \includegraphics[width=\textwidth]{fig_motivating.pdf}
  \caption{One alert, one local model, four ways to split the investigation. The alert's true cause is an exposed admin panel; the planner is Llama-3.1-8B in every row, with the same seed. Each chip is a probe that ran (dashed: a duplicate proposal the controller blocked); red marks the observation that confirms the cause. When the model decides when to stop (first two rows), it reaches the decisive evidence and keeps probing until the budget is gone. A controller-side stop concludes (last two rows), and information-gain selection gets there in three probes. All four rows are taken verbatim from the stopping study's archived journals.}
  \label{fig:motivating}
\end{figure*}

\subsection{Experiment Design}

\subsubsection{Configurations} All configurations use the same model, prompt template, probe catalogue, budgets (10 tool calls, 10 cost units, 12 decisions, 900\,s), duplicate blocking, and verifier. \emph{ReAct-style}: the planner sees the raw observation history only and picks every probe. \emph{Memory-only}: the planner also sees the ledger and still picks every probe. \emph{\hesp{}}: the controller picks the probe (EIG/cost or random), with or without showing its ranking, and with or without the finish guard and the controller-side stop. \Cref{fig:motivating} shows four of these configurations on the same alert. After the stopping study we also ran the frozen controller with the LLM replaced by a script that proposes probes in catalogue order and never finishes. This \emph{LLM-free reference} uses the stopping study's seeds, tasks, budgets, and table (216 episodes); it was not pre-registered and we report it as exploratory.

\subsubsection{Implementation} \hesp{} is about 2{,}900 lines of Python that use only the standard library. Planners are pluggable: one JSON decision protocol serves vLLM and Ollama endpoints and the scripted policies used for LLM-free runs. Each episode gets its own loopback HTTP sandbox, so probes are real HTTP requests with transient failures and state changes, and no episode can observe another. Every study uses Qwen2.5-7B-Instruct (bf16), Qwen2.5-32B-Instruct-AWQ, and Qwen2.5-72B-Instruct-AWQ; the decomposition and stopping studies add Llama-3.1-8B-Instruct (bf16) and Llama-3.1-70B-Instruct-AWQ-INT4. Models are served by vLLM on a single RTX~6000-class GPU at temperature 0.2, with a distinct seed per call. A paired, shuffled, resumable runner executes up to 32 episodes concurrently and refuses to resume if the manifest or the hash of the controller source has changed. After each model, it audits all episodes, archives the journals, and records the archive's SHA-256. The stopping study's 1{,}800 episodes ran in 28 minutes. A suite of 145 unit tests covers the ledger, selectors, blind prompts, guard, controller stop, audit, and the isolation of the table estimator from the environment's generator.

\subsubsection{Evaluation Methodology} Each study's hypotheses, configurations, primary endpoints, and analysis were written down and frozen, together with the source and table hashes, before any of its episodes ran; changes after freezing were appended as errata (Appendix~\ref{app:prereg}). The outcome measure is \emph{verified completion}, the fraction of episodes that end with a verified verdict. Every other ending, including timeouts and malformed outputs, counts as a failure. We report paired per-task differences with percentile task-cluster bootstrap intervals (2{,}000 resamples, 24 clusters). When a study has two primary endpoints, each gets a Bonferroni-adjusted 97.5\% interval; all secondary comparisons are exploratory and uncorrected. We also report descriptive security metrics. Two causes are benign (the authorized scan and the monitoring false positive), and the other six require action. Among episodes that end with a verdict, the \emph{missed-attack} rate is the fraction of cases with an actionable true cause in which the verdict names a benign cause, the \emph{false-escalation} rate is the converse, and the \emph{wrong-cause} rate is the fraction of verdicts that name the wrong cause. The \emph{unresolved} rate is the fraction of all episodes without a verified verdict. Ground truth is the cause in effect at the moment of the verdict. Scripts generate every table and figure from the raw outcome files.
